\cotefile{136}
\scanpage{312}
\mspage{11}
% archive sheet number: 313

la seule connaissance de $C_{\cdot}\in\mathit{Ab}_{\cdot}$, au lieu
et place de $X_{\ast}$ ; il faut, pour \deleted{\unsure{le transport|le transfert}} 
décrire en \uline{termes de} $C_{\cdot}$ \deleted{en} \inserted{certaines des} structures
qu'on peut \inserted{ainsi} décrire sur $F(X_{\ast})$, en termes
de celles de $F$, de tenir compte des
structures supplémentaires convenables
de $C_{\cdot}$, qui peuvent s'expliciter en se
rappelant que $C_{\cdot}$ provient d'un $X_{\ast}$. \deleted{\unsure{Que}}
L'étude systématique de ces structures \deleted{\unsure{qui}}
\deleted{est} \inserted{peut être considéré comme} le contenu essentiel de ces réflexions.

Mais bien sûr, \deleted{si les structures} \inserted{dans les} \unsure{cas} disons
où on a sur $F$ une structure de
$k$-module, i.e. $F:\unsure{\mathscr S|\mathit{Se}|\mathcal{S}s}^{\circ}\unsure{\mathbin{-}|\longrightarrow}\mathit{Ab}_k$
(où $\mathit{Ab}_k$ est la catégorie des $k$-modules),
\deleted{il suffit} \inserted{il est inutile} de se donner de structure
suppl. sur $C^{\unsure{\cdot|\circ|.}}$, car \deleted{\unsure{tous}} grâce au \deleted{fait}
\deleted{\unsure{de compléter par des|d'employer des|de complexe de} $k$-chaînes sur $C_{\cdot}$
(i.e. se donner $k\longrightarrow\operatorname{End}_{\mathit{Ab}_{\cdot}}(C_{\cdot})$ pour
récupérer la structure complète de}
que le $K_{\cdot}$ qui décrit $F$ est
maintenant dans $(\mathit{Ab}_k)_{\cdot}$ i.e. qu'on a

\[
k\longrightarrow\operatorname{End}_{\mathit{Ab}_{\cdot}}(K_{\cdot})
\]

$k$ opère aussi sur $\operatorname{Hom}_{\mathit{Ab}_{\cdot}}(C_{\cdot},K_{\cdot})$,
dont on récupère ainsi la structure
de $k$-module. En \deleted{fait,} il suffit,
au lieu de $C_{\cdot}$, de \unsure{connaitre|connaître}
seulement $\inserted{C_{\unsure{k\ast k|k\ast|k\cdot}}=}C_{\cdot}\otimes_{\mathbb Z} k$ (noté $\mathbb C_{k\ast}(X_{\ast})$),
car on aura

\deleted{\unsure{les relations}}
\[
\left\{\begin{aligned}
F(X_{\ast})&\simeq\operatorname{Hom}_{(\mathit{Ab}_k)_{\cdot}}(C_{k\cdot},K_{\cdot})\\
C_{k\cdot}&=\mathbb C_{k\ast}(X_{\ast})=\mathbb C_{\ast}(X_{\ast})\otimes_{\mathbb Z} k
\end{aligned}\right.
\]

Plus généralement, revenant au cas d'une
$B$ cat. additive générale stable par $\varprojlim$
et $F:(\mathrm{\unsure{Se|S|Ss}}_{\mathrm{f}\unsure{.}})^{\circ}\longrightarrow B$ et $K_{\cdot}\in B_{\cdot}$ se correspondant par la prop. 2, on peut encore
écrire
\[
F(X_{\ast})\simeq\operatorname{Hom}_{\cdot\deleted{\unsure{B}}}^{\unsure{\circ|0|\cdot}}(\mathbb C_{\cdot}(X_{\ast}),K_{\cdot})
\]

en définissant le deuxième mb comme
\uline{objet de $B$}, en définissant d'abord le
bicomplexe de $B$
\[
\begin{aligned}
\mathcal{H}\mathrm{om}^{\cdot}_{\cdot}(C_{\cdot},K_{\cdot})&=
\bigl(\operatorname{Hom}_{B}(C_i,K_j)\bigr)_{i,j}\\
\text{(donc}\qquad \mathcal{H}\mathrm{om}^{i}_{j}(C_{\cdot},K_{\cdot})&=
\operatorname{Hom}_{B}(C_i,K_j)) 
\end{aligned}
\]

-- où l'objet $\mathrm{Hom}(A,M)$, pour $A\in\mathit{Ab}$
et $M\in B$, se définissant par
\[
\operatorname{Hom}_{\uline{B}}(N,\mathrm{Hom}(A,M))\simeq
\operatorname{Hom}_{\mathit{Ab}}(A,\operatorname{Hom}_{B}(N,M))
\]
(le foncteur \deleted{étant \unsure{repr.}} \inserted{du $2^{\text{ème}}$ mb} étant bien
représentable en $N$, pour $A,M$ fixés, grâce
à l'hyp. faite sur $B$ d'être stable par $\varprojlim$).

On définit alors
\[
\operatorname{Hom}_{\cdot}(\mathbb C_{\cdot},\deleted{\unsure{X_{\ast}}},K_{\cdot})
\overset{\mathrm{défn}}{=} Z^0(\mathcal{H}\mathrm{om}^{\cdot}_{\cdot}(\mathbb C_{\cdot},K_{\cdot}))
\]
\inserted{cycles de degré 0 pour degré total (\unsure{cohomol.|contrav.|associé} bien \unsure{ent.}).}

Si enfin $B$ est une cat. $k$-linéaire
($k$ commut.\ \unsure{tjs}) i.e. munie d'un hom d'
anneaux $k\longrightarrow\operatorname{End}(\mathrm{id}_{\uline{B}})$,
